% ===== §8 =====
\section{The Work of the Two Characters}
\label{sec:twocharacters}

\noindent
The four preceding sections passed through trust, shared value, conflict, and the quality of a shared life, and in each the contribution of the present account was carried by two characters, the relational and the generative. It is worth drawing these together, since the functional case rests on them and since it is their conjunction, and not either alone, that gives the account its force. The two characters answer two different questions about any of the matters treated, and the received theories tend to answer each question in a way the account revises.

The relational character answers the question of location, of where the matter is. On the received theories trust is located in the trusting party's estimate of the other, shared value in an overlap of antecedent profiles or a converged consensus, the root of conflict in the disputed matter itself, and the quality of a bond in the satisfaction of individual needs. In each case the matter is located in an individual, in a property of one party, or in a state to be reached. The relational character relocates each matter into the relation between two symbolic worlds. Trust, shared value, the root of conflict, and the quality of a bond are, on the present account, matters generated between the partners and observable only in the relation between their worlds, borne by neither alone. This relocation is the first half of the contribution, and by itself it would yield a relational structuralism, an account on which the matters of intimate life are relational but fixed, given with the structure of the relation and not themselves in movement.

The generative character answers the question of time, of how the matter stands in history. On the received theories the matters of intimate life tend to appear as states, reached when conditions are met and maintained while they hold, trust established, consensus achieved, satisfaction attained. The generative character sets each matter in time as a process and not as a state. Trust is continually regenerated or allowed to thin, shared value is produced anew across a difference and is neither found nor fixed, the easing of conflict generates a growing observation and does not restore a prior peace, and the quality of a bond is a generativity sustained and not a level maintained. This is the second half of the contribution, and by itself, without the relational character, it would yield an account of individual growth, a story of persons developing in time but developing as separate selves, the bond a setting for a change that remains each person's own.

The force of the account lies in the conjunction. The relational character without the generative gives a relation that holds the matters of intimate life but does not move, a structure fixed once the bond is formed. The generative character without the relational gives a movement that belongs to individuals and not to the relation between them, a growth that is not shared. Together they give what the account requires, matters that are borne in the relation between two worlds and that are generated, deepened, and transformed in the historical movement of that relation. Trust, shared value, the easing of conflict, and the quality of a bond are, on the full account, relational matters in generative movement, located between the partners and continually produced in the crossing of their worlds.

\begin{claim}[The division of labor between the two characters]
The relational character answers where a matter of intimate life is located, relocating trust, shared value, the root of conflict, and the quality of a bond from an individual property or a state to be reached into the relation between two symbolic worlds. The generative character answers how the matter stands in time, setting each as a process continually produced and not a state once attained. Neither character suffices alone. The relational without the generative yields a fixed relational structure; the generative without the relational yields a merely individual growth. Their conjunction yields the account's thesis, that the matters of intimate life are relational matters in generative movement, borne between two worlds and produced in the history of their crossing.
\end{claim}

\begin{table}[h]
\centering
\small
\setlength{\extrarowheight}{3pt}
\begin{tabular}{@{}p{2.5cm} p{3.5cm} p{3.3cm} p{3.3cm}@{}}
\arrayrulecolor{gold}\specialrule{0.7pt}{0pt}{2pt}
{\color{rosedeep}\bfseries Function} &
{\color{rosedeep}\bfseries Received rendering} &
{\color{rosedeep}\bfseries Relational character adds} &
{\color{rosedeep}\bfseries Generative character adds} \\
\specialrule{0.5pt}{2pt}{3pt}
Trust &
an estimate, by one party, of a property or probability of the other &
a shared reality borne between the two, extended before the evidence a wager would require &
continually regenerated in the crossing, thinning where it is no longer renewed \\[3pt]
Shared value &
a pre-existing overlap found, or a consensus in which difference is overcome &
a matter held in common across a difference the two do not dissolve &
produced anew in the relation, arising where neither world held it before \\[3pt]
Conflict &
a disagreement about the disputed matter, corrosive or reparable in its patterns &
a collision of two unrecognized renderings, moved from content to the level of worlds &
each crossing generating a growing observation of the difference, a resource for the next \\[3pt]
Quality of a shared life &
a level of satisfaction of the individual's needs, higher or lower &
the richness of what the two worlds generate between them, borne by neither alone &
a generativity sustained, a spiral that returns and finds the shared life grown \\
\specialrule{0.7pt}{3pt}{0pt}
\end{tabular}
\caption{The functional case in summary. For each of the four functions, the received theories render the matter as an individual property or a state to be reached; the relational character relocates it into the relation between two worlds, and the generative character sets it in historical movement. The account's thesis is the conjunction of the two additions, that the matters of intimate life are relational matters in generative movement.}
\label{tab:functions}
\end{table}

This completes the functional case. The account has set the potential work of relational understanding beside the established theories that treat the same matters, and it has located its contribution in two characters whose conjunction supplies what those theories, rendering the matters of intimate life as individual properties or states, leave unsupplied. A reader persuaded to this point might conclude that the case for relational understanding in intimacy is made. The next part argues that it is not, and that a case built on function alone, however strong, leaves open a question that function cannot answer.
